\begin{tcolorbox}[breakable,title=Prompt Details for Academia]
\textbf{\textcolor{blue}{System Prompt}}\\ 
You can use actions to help people solve problems.\\
\begin{tikzpicture}
\draw[dashed] (0,0) -- (\linewidth,0); % 绘制一条宽度为\linewidth的虚线
\end{tikzpicture}
\textbf{\textcolor{blue}{Instruction}}\\
We detail name, description, input(parameters) and output(returns) of each action as follows:\\
Name: loadPaperNet()\\
Description: Load PaperNet. In this net, nodes are papers and edges are citation relationships between papers.\\
\\
Name: loadAuthorNet()\\
Description: Load AuthorNet. In this net, nodes are authors and edges are collaboration relationships between authors.\\
\\
Name: neighbourCheck(graph, node)\\
Description: List the first-order neighbors connect to the node. In paperNet, neigbours are cited papers of the paper. In authorNet, neigbours are collaborators of the author.\\
Parameters:\\
- graph (Type: string, Enum: [PaperNet, AuthorNet]): The name of the graph to check\\
- node (Type: string): The node for which neighbors will be listed\\
Returns:\\
- neighbors (Type: array)\\
\\
Name: paperNodeCheck(node)\\
Description: Return detailed attribute information of a specified paper in PaperNet\\
Parameters:\\
- node (Type: string): Name of the paper.\\
Returns:\\
- authors : The authors of the paper\\
- year : The puslished year of the paper\\
- venue : The published venue of the paper\\
- n\_citation : The number of citations of the paper\\
- keywords : The keywords of the paper\\
- doc\_type : The document type of the paper\\
\\
Name: authorNodeCheck(node)\\
Description: Return detailed attribute information of a specified author in AuthorNet\\
Parameters:\\
- node (Type: string): name of the author.\\
Returns:\\
- name : The name of the author\\
- org : The organization of the author\\
\\
Name: authorEdgeCheck(node1, node2)\\
Description: Return detailed attribute information of the edge between two specified nodes in a AuthorNet.\\
Parameters:\\
- node1 (Type: string): The first node of the edge\\
- node2 (Type: string): The second node of the edge\\
Returns:\\
- papers : All papers that the two authors have co-authored\\
\\
Name: paperEdgeCheck(node1, node2)\\
Description: Return detailed attribute information of the edge between two specified nodes in a PaperNet.\\
Parameters:\\
- node1 (Type: string): The first node of the edge\\
- node2 (Type: string): The second node of the edge\\
Returns:\\
None\\
\\
Name: check\_valid\_actions()\\
Description: Get supported actions for current tool.\\
Returns:\\
- actions (Type: array): Supported actions for current tool.\\
\\
Name: finish(answer)\\
Description: Return an answer and finish the task\\
Parameters:\\
- answer (Type: [`string', `number', `array']): The answer to be returned\\
\\
\\
If you are finished, you will call ``finish'' action\\
Please refer to the format of examples below to solve the requested goal. Your response must be in the format of ``Action: [your action] with Action Input: [your action input]''\\
\begin{tikzpicture}
\draw[dashed] (0,0) -- (\linewidth,0); % 绘制一条宽度为\linewidth的虚线
\end{tikzpicture}
\textbf{\textcolor{blue}{Examples}}\\
Goal: When was the paper Learning the Principle of Least Action with Reinforcement Learning. published?\\
\\
Action: loadPaperNet with Action Input: \{\}
\\
Observation: PaperNet is loaded.\\
Action: paperNodeCheck with Action Input: \{``node'':``Learning the Principle of Least Action with Reinforcement Learning.''\}\\
Observation: \{`year': 2021, `venue': `AAAI Spring Symposium - MLPS', `n\_citation': 0, `keywords': [], `doc\_type': `Conference'\}\\
Action: finish with Action Input: \{``answer'': ``2021"\}\\
Observation: 2021
\end{tcolorbox}
\begin{figure}[!h]
    \vspace{-0.08cm}
    \caption{Prompt Details for Academia in Tool-Query Environments.}
    \label{fig:academia_prompt}
\end{figure}